\section{Before Vatican I: Primacy, Revolution, and the Theologies of Authority}

Vatican I cannot be understood if the papacy is made either an invention of 1870 or a timeless office whose exercise never changed. Catholic doctrine confesses a Petrine ministry instituted by Christ and continuing in the bishops of Rome. History records changing legal forms, contested claims, uneven communication, cooperation with councils and patriarchates, resistance by rulers and bishops, and a nineteenth-century concentration of Roman authority. Development concerns the articulation and exercise of a given office; it does not require pretending that every century looked like the next.

\subsection{The inheritance before the modern crisis}

The scriptural case received by the Catholic tradition joins Matthew 16:16--19, Luke 22:31--32, John 21:15--17, and the Petrine role in Acts. Early Roman interventions such as \emph{1 Clement}, the Roman succession list and “pre-eminent” Church in Irenaeus (\emph{Against Heresies} III.3), appeals to Rome, and the doctrinal role of Leo I at Chalcedon are real evidence. They do not form a modern canon-law manual. First-millennium primacy operated through communion, synods, reception, imperial structures, patriarchal rights, and contested jurisdiction. Catholic development must preserve that evidence rather than make the mature formula appear verbatim in every age.

Medieval conflict made the institutional claims sharper. The Gregorian reform defended ecclesial freedom against lay investiture; canonists developed appellate and plenary jurisdiction; Boniface VIII's \emph{Unam sanctam} (1302) joined spiritual claims to a particular medieval political idiom. The Great Western Schism then made the relation of pope and council an existential problem. Constance ended the competing obediences, while later Catholic reception did not accept a standing conciliar superiority over a legitimate pope. Florence (1439) taught a universal Roman primacy in the setting of an attempted East--West reunion. Trent (1545--1563), by contrast, did not define papal infallibility; its decrees were papally confirmed and implemented through an increasingly Roman confessional system.

This long inheritance establishes two facts at once. The Petrine office is not a response invented against the French Revolution. The juridical centralization and popular papal devotion of the nineteenth century are not identical with every earlier exercise of the office.

\subsection{Gallicanism, Febronianism, and the national Church}

The proximate adversary of \emph{Pastor aeternus} was not simply Protestantism. It was a Catholic family of limitations on Roman power. The Four Gallican Articles adopted by the French clergy in 1682 denied papal authority over temporal rulers, subordinated papal judgments to the received canons and customs, honored the authority of general councils in the conciliarist sense, and held papal judgments reformable unless the Church consented. Different Gallicans emphasized different elements; not every defender of local episcopal rights held the whole package.

Febronianism in the German lands and Josephinism in the Habsburg monarchy joined episcopal or national-Church theories to state administration. Crown nomination, censorship of Roman acts, control of seminaries and religious houses, and suppression of “unproductive” monasteries showed that decentralization from Rome could mean not the freedom of local Churches but their subjection to ministers. \emph{Pastor aeternus}'s insistence that the pope communicate freely with pastors and faithful, and that his jurisdiction is not mediated by civil permission, answers this political experience as well as an abstract ecclesiology.

\begin{threestable}{Current}{Theological or political claim}{Pressure toward Vatican I}
Conciliarism and Gallicanism & Final doctrinal judgment required a council or subsequent ecclesial consent; inherited canons and national usages limited papal exercise. & Rome sought to define where the Church's final visible judgment resides and to exclude a ratification veto.\\
Febronian and Josephinist reform & Episcopal or national structures could be reorganized through enlightened state power. & The papacy appeared as a transnational defense of episcopal and ecclesial freedom.\\
Revolutionary sovereignty & Church property, corporations, orders, and clerical oaths were made subject to the nation. & Catholics sought an authority not created by the state and capable of unifying dispersed believers.\\
Ultramontanism & Catholics looked “beyond the mountains” to Rome for doctrine, devotion, discipline, and resistance to national control. & Primacy became a lived mass identity as well as a theological thesis.\\
Liberal Catholicism & Constitutional freedoms and a free Church could be reconciled; some advocates also pressed wider theological accommodation. & It exposed divisions over whether the modern state was chiefly threat, opportunity, or both.\\
\end{threestable}

\subsection{1789 and the destruction of an order}

The French Revolution did not merely change a ministry. Nationalization of Church property, suppression of religious vows and houses, and the Civil Constitution of the Clergy (1790) attempted to reconstruct dioceses and select bishops through the new sovereignty. The required oath divided clergy and faithful. Revolutionary violence, dechristianization, and the execution or exile of clergy gave later Catholic memory a martyrological as well as political cast. Napoleon's concordat with Pius VII restored public worship but also demonstrated the ruler's capacity to reorganize episcopal structures; the pope's captivity gave personal form to the conflict.

After 1815 the restored alliance of throne and altar did not restore the old world. Rail, telegraph, cheap print, pilgrimage, new congregations, and popular devotions made Rome more accessible. The papacy lost older temporal protections while gaining unprecedented moral visibility. Ultramontanism was therefore not only a theory imposed downward. Lay movements, journalists, new religious institutes, and clergy often desired Roman uniformity as liberation from hostile governments and local elites.

\subsection{Reason under pressure: rationalism and its Catholic reactions}

The first constitution of Vatican I was not about the pope. \emph{Dei Filius} answered a wider crisis of knowledge. Enlightenment rationalism could restrict truth to autonomous reason and deny supernatural revelation or miracle. Materialism and pantheism threatened the Creator--creature distinction. Historical criticism raised new questions about Scripture and dogma. Yet the Catholic reaction could also fail: fideism and traditionalism sometimes so distrusted fallen reason that they denied its natural capacity to know God or made all fundamental knowledge depend on an original social revelation.

Nineteenth-century Catholic theology was not a two-party duel between modernists and scholastics. Louis de Bonald and Joseph de Maistre defended authority through traditionalist accounts of society. Félicité de Lamennais moved from ultramontane defense toward a liberal Catholic program and finally rupture. Georg Hermes and Anton Günther sought rational reconstructions that Roman judgments found inadequate. The Roman School—Giovanni Perrone, Carlo Passaglia, Clemens Schrader, Johann Baptist Franzelin—developed the relation of Scripture, Tradition, dogma, and living magisterium. Johann Adam Möhler described the Church as a Spirit-filled organism. John Henry Newman analyzed development and the active role of the whole faithful in reception without turning doctrine into majority vote.

The council's answer was a double refusal. Reason can know God with certainty from created things; faith is nevertheless a supernatural virtue and response to divine Revelation. Mysteries exceed reason but do not contradict it. Sciences have legitimate methods, yet cannot contradict a truth truly known from God. John Paul II later described \emph{Dei Filius} as censuring both distrust of reason and rationalist excess (\emph{Fides et ratio} 52--53). This balance is indispensable: Vatican I was not simply the enthronement of will over reason.

\subsection{Pius IX and the tightening siege}

Pius IX's long pontificate began in 1846 amid reforming expectations. The revolutions of 1848, the murder of his minister Pellegrino Rossi, the pope's flight to Gaeta, and restoration through French arms hardened the political horizon. Italian unification steadily reduced the Papal States. By the time the council met, Rome depended on a French garrison and faced imminent incorporation into the Kingdom of Italy.

Two preconciliar acts shaped expectations. In \emph{Ineffabilis Deus} (1854), Pius IX defined the Immaculate Conception after consulting the bishops. The act supplied a recent paradigm of a papal definition: personally issued, but preceded by inquiry into the Church's faith. The \emph{Syllabus of Errors} (1864), appended to \emph{Quanta cura}, collected propositions from earlier acts. Torn from their originating contexts, its terse negatives were easily read as a comprehensive rejection of every modern freedom. The episode already displayed a recurrent problem: centralized propositions travel farther than their interpretive setting.

In 1867 Pius IX announced his intention to convoke a council; the bull \emph{Aeterni Patris} summoned it for 8 December 1869. The preparatory machinery produced schemas on faith, the Church, bishops, clergy, religious life, missions, and discipline. Infallibility was not the council's only contemplated subject. It became urgent through petitions and counter-petitions, organized ultramontane advocacy, public journalism, and the belief that the shrinking temporal sovereignty made a spiritual definition providentially necessary.

\judgmentbox{Theology and siege cannot be separated, but neither can one reduce doctrine to siege psychology. Gallican consent theory, state control of national Churches, rationalism, and the need for final doctrinal judgment were real theological problems. The loss of papal territory, revolutionary memory, new communications, and Pius IX's personality shaped the timing, rhetoric, and compression of the answer.}
